\section{结果讨论与政策含义}
\subsection{结果讨论}
结合实证章节结果可以看出，我国碳排放变化并非由单一变量决定，而是规模因素与结构因素共同作用的结果。扩展STIRPAT估计显示，在2005—2022年小样本条件下，仅煤炭占比在95\%置信水平下保持显著正向影响，是统计上最稳健的驱动因子。人均GDP、能源强度、城镇化率与第二产业比重等变量点估计方向虽与经济直觉一致，但区间较宽且跨越零，提示其作用强度识别仍受自由度约束。

与此同时，城镇化率与技术进步变量的估计结果并不稳健。城镇化系数虽为正，但置信区间较宽，提示其作用方向受模型设定和样本长度影响较大；技术进步变量在 STIRPAT 弹性估计中呈现极微弱的减排倾向，但在中介效应分析中却表现出微弱的正向关联。这种系数符号的不确定性，反映了在小样本约束下技术指标对碳排放的影响存在‘结构减排’与‘技术反弹’的复杂博弈， 同时也提示现有专利数据可能存在传导时滞或未能完全覆盖生产工艺的深度绿色化。因此，未来政策应更关注低碳技术的实质转化率而非单纯的数量扩张。

中介路径结果进一步支持协同机制判断。能源结构不仅直接影响碳排放，还可能通过技术进步形成间接作用链。换言之，能源替代并不是孤立政策动作，而是会牵引技术改造、设备更新和产业重构的系统过程。因此，若政策仅停留在末端治理，往往难以改变高碳路径惯性；只有将能源、产业、技术和治理政策协同推进，才能在中长期形成可持续的减排斜率。

预测比较结果揭示了小样本条件下的严重过拟合现象：在统一可比窗口下 AR2 模型表现最优，朴素模型次之；而 XGBoost 与随机森林等复杂机器学习模型因严重过拟合，误差反而比朴素模型高出数倍。这说明在国家级小样本时序条件下，盲目追求算法复杂度不仅无益，反而会引入较大的预测偏差。模型适配性（结构与数据规模的匹配度）才是决定外推可靠性的关键。

情景模拟进一步表明，不同政策强度下的路径差异主要体现为峰值高度与峰后下降速度差异，而不仅是名义达峰时点差异。低碳与深度脱碳情景均可显著压低2035年前后的排放水平，说明政策力度的提升能够有效降低未来减排成本并改善达峰后转型空间。由此可见，政策目标应从“是否达峰”进一步拓展为“以多低峰值达峰、以多稳斜率下降”。

\subsection{政策含义}
基于上述结果，国家层面政策设计可围绕三条主线展开。第一，将能源结构调整作为绝对主抓手，稳步降低煤炭依赖并提高新能源消纳能力。第二，推动第二产业绿色化与高端化，优先在高耗能行业实施节能改造与流程再造。第三，强化技术创新的减排转化效率，政策重点应从``投入规模扩张''转向``低碳技术的产业化应用''。

从政策时序看，应划分为``控强度—压峰值—稳下降''三个阶段：近期重点控制能源强度和高碳增量；中期重点降低煤炭占比以压低峰值平台；远期通过技术扩散和制度协同提高峰后下降稳定性。该分阶段思路将本文的统计量化结果转化为了可执行的政策节奏。

\subsection{稳健性与解释边界}
本文结论适用于国家总量层面，2023—2024年数据仅用于预测延展，不参与机制估计，以保证推断口径一致。

总体而言，本文的核心启示是：在保持经济稳定运行的约束下，我国碳达峰路径优化的关键不在单一指标改善，而在能源结构、产业结构、技术进步和治理之间实现协同发力。只有将“压峰值”与“稳下降”并重，才能提升达峰后的持续减排能力和长期转型质量。
